\chapter*{概论}
\addcontentsline{toc}{chapter}{概论}

这本讲义围绕三条相对独立、但在方法上彼此贯通的主线组织：代数与分析工具、微分几何、概率与数理统计。共同原则不是把不同分支强行归入一个口号，而是始终先确定对象、空间、映射和定义域，再从一般结构推出可计算公式，最后才进入低维表示、可解模型或统计特例。

第一部分从群作用与表示论开始，依次建立结合代数与 Clifford/Spin 结构、Hilbert 空间与广义本征展开、闭算子谱理论、分布与 Fourier 分析、Fredholm 理论以及非线性演化。这样安排使群表示、旋量、谱测度、Green 算子和可积系统不再是孤立技巧，而都能追溯到各自的母结构。第二部分从切空间、外微分和流出发，建立标架、一般向量丛联络、挠率与曲率、Levi--Civita 几何、Hodge 理论、Frobenius 可积性、几何演化以及任意余维子流形的 Gauss--Codazzi--Ricci 方程。第三部分从概率空间和随机样本开始，沿极限定理、充分性与完备性、似然/Fisher 信息、Bayes 推断、假设检验，最终进入一般线性模型、ANOVA 与回归。

三部分都把“精确一般理论 → 参数或结构的具体化 → 受控近似 → 特例”作为单向推导顺序。正文负责给出数学图像、定义与关键结论；较长的证明和代数计算集中在推导环境中，使主线可以独立阅读，而关键结论仍保持可审计性。
